\documentclass[11pt]{article}
\usepackage{amsmath,amssymb,amsthm}
\usepackage[margin=1in]{geometry}
\usepackage{booktabs}
\usepackage[colorlinks=true,linkcolor=blue,citecolor=blue,urlcolor=blue]{hyperref}
\newtheorem{proposition}{Proposition}
\newtheorem{lemma}{Lemma}
\theoremstyle{remark}
\newtheorem{remark}{Remark}

\title{Selection conventions and trial factor in a proposed closed form\\ for the fine-structure constant}
\author{universalunifyingmath project}
\date{Draft, 30 October 2026}

\begin{document}
\maketitle

\begin{abstract}
Wissner-Gross [1] proposes
$\alpha_{\rm cand}=(79+48\sqrt2+42\sqrt3+3\sqrt6)/31104$, with $\alpha_{\rm cand}^{-1}=137.035\,999\,189\,566\ldots$, as a parameter-free candidate for $\alpha(0)$. It is derived from an axiom about maximally entangled qudit pairs. We report four results.
(i) An independent operator-level reconstruction reproduces the $172$ records, the sector sizes, and the value exactly.
(ii) The two ``commuting exchanges'' that sort the records are a single involution on all product records. The weight law is therefore a function of exchange parity and slot class alone.
(iii) The selection of the qutrit--qubit link by an ``acting central charge'' depends on a conjugation convention for the far end of a link, which the axiom does not fix. Under the conjugate convention a second configuration satisfies every hypothesis of the uniqueness theorem. We verify this in Lean~4 against the author's own definitions.
(iv) In a declared space of $46\,080$ rule variants that contains every alternative listed in [1], the expected number of chance values within $1\sigma$ of CODATA~2022 is $4.4\times10^{-6}$. We discuss what this does and does not imply.
\end{abstract}

\section{Setting}
We follow the notation of [1]. The object is a qutrit shared by three qubit links. Two of them are read, and the qutrit is carried as a copy per link. The four \emph{slots} $Q_1,q_1,Q_2,q_2$ carry the letters: Gell-Mann matrices at $Q_i$ and Pauli matrices at $q_i$. The records are
\begin{itemize}
\item the $169$ two-slot products, recombined into eigenvectors of the slot exchange $E$ (induced by $(Q_1,q_1,Q_2,q_2)\mapsto(Q_2,q_2,Q_1,q_1)$);
\item three commutator records $[\sigma_\mu,\sigma_\nu]^{(q_1)}-[\sigma_\mu,\sigma_\nu]^{(q_2)}$.
\end{itemize}
The value is $N^{-1}\sum_r\sqrt{k_r}$ with $N=24\cdot36^2$.

\section{Reproduction}
We built the records as explicit $36\times36$ operators on $\mathbb C^3\otimes\mathbb C^2\otimes\mathbb C^3\otimes\mathbb C^2$ and checked the following:
\begin{itemize}
\item the slot-pair sizes are $(24,24,64,24,24,9)$;
\item $E$ fixes $11$ products and pairs $158$;
\item the $172$ records are Hilbert--Schmidt orthogonal;
\item the sector sizes are $79/48/42/3$.
\end{itemize}
Exact summation gives the closed form, and its inverse agrees with [1] to $60$ digits. The author's Lean file builds under Lean~4.35.0 with axioms \texttt{propext}, \texttt{Quot.sound} and \texttt{Classical.choice} only.

\section{The two exchanges coincide}
\begin{proposition}\label{prop:class}
On the span of the product records, the exchange of reads $X$ and the exchange of links $Y$ of [1] act by the same involution $E$. Let $e_r\in\{\pm1\}$ be the $E$-eigenvalue of record $r$. The squared weight of [1, Eq.~(5)] is
\[
k_r=\begin{cases}1,& e_r=-1 \text{ (product)},\\ 3,& e_r=+1,\ \text{both letters at one end},\\ 2,& e_r=+1,\ \text{one letter at each end},\\ 6,& \text{commutator record}.\end{cases}
\]
\end{proposition}
\begin{proof}
On slots, both exchanges are the permutation $Q_1\leftrightarrow Q_2$, $q_1\leftrightarrow q_2$. They differ only in the set of records each is declared to act on: same-end records for $X$, mixed records for $Y$, and commutator records for both [1, Sec.~S7]. Substituting into $k_r=3^{[X_r=-1]}2^{[Y_r=-1]}$ gives the table. In Lean this is \texttt{weight\_is\_class\_function}, proved by kernel evaluation on the author's definitions.
\end{proof}
Proposition~\ref{prop:class} isolates the content of the weight law in two assignments, $\sqrt3$ to even same-end records and $\sqrt2$ to even mixed records. In [1] these follow from wording clauses (``a read is an act at the end at which it is taken''; ``the copies of a shared qudit are one qudit''). Rewording them gives $1/\text{value}\in[112.4,146.5]$ [1, Table S5].

\section{Convention dependence of the selection}
The identification algebra of a $(d_1,d_2)$ link in [1] is $s(u(d_1)\oplus u(d_2))$, traceless on $\mathbb C^{d_1}\oplus\mathbb C^{d_2}$. It acts on $\mathbb C^{d_1}\otimes\mathbb C^{d_2}$ by $\xi_1\otimes 1+1\otimes\xi_2$. The pair state, however, is invariant under $U\otimes\bar U$, so the far end transforms contragrediently.

\begin{proposition}\label{prop:conv}
Let $c=(ix\mathbb 1_{d_1},iy\mathbb 1_{d_2})$ be central with $d_1x+d_2y=0$, $x\neq0$.
\begin{enumerate}
\item[(A)] If the far block acts by $\xi_2$, then $c$ acts on the link by $i(x+y)$. This is nonzero iff $d_1\neq d_2$.
\item[(A$^*$)] If the far block acts by $\bar\xi_2$, then $c$ acts by $i(x-y)=ix(d_1+d_2)/d_2\neq0$ for every type.
\item[(B)] If instead the phase constraint is imposed on the matched directions, $x+y=0$, the centre acts trivially for every type.
\end{enumerate}
Among readable types $\{(2,2),(2,3)\}$, reading A selects $(2,3)$ alone, A$^*$ selects both, and B selects neither. Under A the link has hypercharge $1/6$, the quantum numbers of $Q_L$. Under A$^*$ it has $5/6$, those of the $(3,2)$ block of the $SU(5)$ adjoint.
\end{proposition}
\begin{proof}
Direct substitution of $y=-d_1x/d_2$ in (A) and (A$^*$). In (B) the action is identically zero.
\end{proof}

\begin{remark}
The trace condition defining the $u(1)$ is not invariant under $\xi_2\mapsto\bar\xi_2$: since $\operatorname{tr}\bar\xi_2=-\operatorname{tr}\xi_2$, it becomes $\operatorname{tr}\xi_1-\operatorname{tr}\xi_2=0$. The two readings therefore define different subgroups of $U(d_1)\times U(d_2)$. Both give a four-dimensional stabilizer of the $(2,3)$ pair state, so the stabilizer does not discriminate between them.
\end{remark}

In Lean (appended to the unmodified file of [1]) we prove \texttt{uniqueness\_fails\_under\_conjugate\_convention}. With the acting-charge predicate replaced by its A$^*$ form, the configuration ``shared qubit, two read and one unread $(2,2)$ link'' satisfies all hypotheses of the upstream uniqueness theorem and differs from the object. Theorem~1 of [1] is thus correct under A. Its scope depends on a convention whose justification in [1, Sec.~S4] uses the direct sum $\mathbb C^{d_1}\oplus\mathbb C^{d_2}$, a space the axiom does not supply.

\section{Trial factor}
All $13$ closed-form alternatives in [1, Table S5] are reproduced from the record classes. We then declared a rule space that contains them:
\begin{itemize}
\item an independent squared factor $k\in\{1,2,3,6\}$ for each of the four product classes;
\item a commutator factor in $\{1,2,3,4,6,9\}$ and count in $\{0,3,6\}$;
\item amplitude or Born weighting;
\item five normalizations $N\in\{24\cdot36^2,36^3,24^3,24\cdot35^2,35^3\}$.
\end{itemize}
This gives $46\,080$ rules and $15\,639$ distinct values. The rule of [1] is the closest to CODATA~2022, at $0.598\sigma$. The next closest is at $1.3\times10^5\sigma$. Near $1/\alpha$ the distinct values have density $103.9$ per unit, so the expected number in the $\pm1\sigma$ window ($4.2\times10^{-8}$) is $4.4\times10^{-6}$.

Separately, among the $1.63\times10^9$ elements $(a+b\sqrt2+c\sqrt3+d\sqrt6)/31104$ with $0\le a,b,c,d\le200$, exactly one lies within $1\sigma$: $(79,48,42,3)$. The expected number is $0.024$.

These numbers bound coincidence only relative to a stated search space. If the rules were fixed without reference to $\alpha$, the agreement is improbable at the $10^{-6}$ level. If they were selected from an effective space of $10^{9}$--$10^{10}$ variants, it is not. Variants of this kind include wording clauses, configurations, normalizations, and conventions such as that of Proposition~\ref{prop:conv}. [1] does not document the order in which rules and value were fixed, so the material cannot decide between the two.

\section{Testability}
The candidate deviates from the measurements as follows: $+0.60\sigma$ (CODATA 2022), $+1.57\sigma$ (electron $g-2$), $-1.49\sigma$ (Rb), $+5.32\sigma$ (Cs). Consistent determinations with $u\lesssim2.5\times10^{-9}$ would test it at $5\sigma$, if centred near the present CODATA value. The identification also requires a map from the hypercharge readout to $\alpha_{\rm em}=\alpha_Y\cos^2\theta_W$, which [1] defers.

\section*{Reproducibility}
Code, outputs and Lean files are in the project repository, under \texttt{computation/chain-theory/alpha/} and \texttt{proofs/lean/chain-theory/}. Upstream sources are pinned to commit \texttt{a28db89} of \texttt{github.com/alexwg/chain-theory}.

\begin{thebibliography}{9}
\bibitem{wg} A.~D.~Wissner-Gross, \emph{A Parameter-Free Candidate for the Fine-Structure Constant from Entangled Qudits}, preprint v1.0.0 (2026), with Supplemental Material and \texttt{ChainAlpha.lean}.
\bibitem{codata} P.~J.~Mohr, D.~B.~Newell, B.~N.~Taylor, E.~Tiesinga, Rev.\ Mod.\ Phys.\ \textbf{97}, 025002 (2025).
\bibitem{fan} X.~Fan, T.~G.~Myers, B.~A.~D.~Sukra, G.~Gabrielse, Phys.\ Rev.\ Lett.\ \textbf{130}, 071801 (2023).
\bibitem{morel} L.~Morel, Z.~Yao, P.~Clad\'e, S.~Guellati-Kh\'elifa, Nature \textbf{588}, 61 (2020).
\bibitem{parker} R.~H.~Parker, C.~Yu, W.~Zhong, B.~Estey, H.~M\"uller, Science \textbf{360}, 191 (2018).
\end{thebibliography}
\end{document}
